\section{Introduction}
\label{sec:introduction}

In the maize departments studied here, July and August 2003 brought \HeatTwoThousandThree{} degree-days above 32\,$^{\circ}$C on average, \HeatRatio{} times more than any other summer since 2000 (\cref{fig:projections}). Maize suffered, but not all maize equally: in the same department, fields with irrigation and fields without it did not lose the same share of their harvest. This monograph asks how large that difference is, how it moves with heat, and how sure we can be about it.

The question matters because irrigation is expensive and because its water is scarcest exactly when it is most valuable. In hot, dry summers, French prefectures issue orders that restrict water withdrawals and can limit irrigation (\cref{sec:data}). Knowing how the yield advantage of irrigated maize varies with summer heat is a first step towards valuing that water.

\subsection*{What we already know}

Extreme heat lowers the yields of the major field crops, in France as elsewhere \parencite{schlenker2009,gammans2017}. Studies abroad find that irrigation weakens this association. In the United States, yield declines on hot days are much weaker in irrigated areas than under rainfed conditions \parencite{schauberger2017}, and the benefit of irrigation depends on growing-season climate \parencite{li2018}. In India, irrigated wheat shows about a quarter of the heat sensitivity of rainfed wheat \parencite{zaveri2019}. Field trials in Kansas point the same way \parencite{tack2017}.

For France, the evidence is indirect. \textcite{hawkins2013} infer a role for irrigation from the declining sensitivity of maize yields to rainfall. \textcite{vandervelde2010} simulate irrigated and rainfed maize with a crop model during the 2003 heatwave and weight the simulations by regional irrigated shares. I am not aware of a statistical estimate built on observed irrigated and rainfed yields.

\subsection*{What this monograph does}

\paragraph{It compares irrigated and rainfed maize side by side.} The French annual agricultural statistics (\emph{Statistique agricole annuelle}) report grain maize yields separately for irrigated and rainfed land in each department. I located an older release that extends this split back to 2000 and checked it against the current one, which gives 26 harvests. The outcome is the log yield gap between the two modes in the same department and year. Shocks that hit both modes proportionally in a given year, for instance a disease outbreak or a late frost if it hits both modes alike, cancel out in the gap. Stable differences between irrigated and rainfed land, such as soils or farm types, are absorbed by department effects. What is left is how the advantage of irrigated maize moves with summer heat.

\paragraph{It measures heat where the maize grows.} Degree-days above 32\,$^{\circ}$C are computed for each 8~km grid cell and each day, then averaged over each department with weights given by the area of maize parcels. The choice matters. In the \ReliefCount{} departments with uplands, a department-wide average understates the heat that maize experienced in July--August 2003 by \ReliefGapSmallMin{} to \ReliefGapMax{} degree-days (\cref{tab:relief}; \cref{fig:map} shows where they lie), and with that measure the main result is no longer significant (\cref{sec:maize}).

\paragraph{It checks whether its own tests can be trusted.} This is where the project took its most important turn. For wheat, heat at the end of the cycle was associated with \WheatPct\,\% lower yields per degree-day. Two of the three inferences planned in advance called this significant: $p = \WheatPdep$ with standard errors clustered by department. The third, a bootstrap by year, did not ($p = \WheatPwcr$). Which one to believe? I asked a simple question: if heat had no effect at all, how often would each test still say it does? To find out, I replaced each year's heat with the heat of another year, keeping its real geography, and re-ran the tests hundreds of times. Clustering by department said ``significant'' \WheatSizeDep\,\% of the time, instead of the nominal 5\,\%. A hot summer is hot everywhere, so 83 departments are not 83 independent experiments, and the usual standard errors behave as if they were.

The problem is not specific to wheat or to this project. Across the five size checks run here, on three panels (\cref{tab:sizes}), standard errors clustered by department and \textcite{conley1999} standard errors with a 200~km cut-off (uniform kernel, no serial component, as implemented here) reject a true null in \UsualSizeMin\,\% to \UsualSizeMax\,\% of cases. A wild cluster bootstrap with clusters by year and Webb weights \parencite{cameron2008,webb2023} stays close to 5\,\%, and it is the only inference on which the conclusions below rely. This size check was not in the original plan; it was added as a diagnostic, and it is now applied to every main test. \Cref{sec:inference} explains the check so that others can run it on their own data; the code is a stand-alone module (\cref{app:wcr}).

\paragraph{It fixed its rules in advance, and says where it did not.} Each estimation phase began with a written specification whose SHA-256 hash was recorded in the project's decision log before any estimate of that phase was computed. The log is published with the code, but it was kept by the author alone, and the specifications, the log and the results are published together. The hashes therefore show only that the files have not changed since publication; that each specification came before its estimates, and that each check added later was declared before it was computed, is the author's statement. For maize, the decision rule was set at the same time: a hypothesis counts as established only if its one-sided bootstrap $p$-value is below 0.05 \emph{and} the measured size of the test is at most 10\,\%. The maize specification was not written blind, since an exploratory analysis of 2010--2025 already existed; this is why the untouched years 2000--2009 serve as a replication. Some analyses were added after results were known: the size checks themselves, extra robustness checks for wheat, and the comparison of the forecasts with the department mean (B0), a benchmark defined in the specification but not used as its criterion. They are reported as such, and every departure from the plans is listed in \cref{app:deviations}.

\subsection*{What I find}

In the 17 departments with the best quality status (the 2025 status, applied to the whole period), each degree-day above 32\,$^{\circ}$C in July--August is associated with a \QonePct\,\% wider gap in favour of irrigated maize, that is, a \QonePct\,\% higher ratio of irrigated to rainfed yield: at the median, an advantage of \GapMedianPct\,\% would become \GapMedianAfterPct\,\% (\cref{fig:maize-resid}; $p = \QoneP$; the test's measured size is \QoneSizeWCR\,\% and \PowerCheckSize\,\% in two placebo runs). Each degree-day is associated with about \QoneRainfedPct\,\% lower rainfed yields and \QoneIrrigatedPct\,\% lower irrigated yields. On these point estimates, rainfed maize is thus about \QoneRatio{} times more sensitive to heat (the size of these two tests, measured afterwards, is \HtwobSize\,\% and \HtwocSize\,\%).

This meets the rule fixed in advance, but the evidence is weaker than the $p$-value suggests. \ExploSharePct\,\% of the department-years belong to 2010--2025, the period whose exploratory analysis suggested the hypothesis; the only new evidence is the decade 2000--2009. Its magnitude is uncertain: the 90\,\% interval obtained by inverting the bootstrap runs from \CiNinetyLow{} to \CiNinetyHigh\,\% per degree-day. With the maize mask, the estimate is \MaskRecent{} on 2010--2025 and \MaskEarly{} on 2000--2009; the difference comes from 2003, since without that summer the earlier decade gives a larger estimate (\cref{tab:maize-measure}). The exploratory analysis, with department-wide weather, had given \AreaRecent{} on 2010--2025, and that measure gives $\AreaEarly$ on 2000--2009.

Eight checks shape how far this result can be trusted.
\begin{itemize}
  \item \textbf{New years.} On 2000--2009 alone, a decade no earlier analysis had touched, the estimate keeps its sign with the maize mask, at about half the 2010--2025 estimate (\QoneReplPct\,\% against \QoneRecentPct\,\%, one-sided $p = \QoneReplP$; the specification set no threshold for the replication). With department-wide weather it does not.
  \item \textbf{An older map of maize.} With a 2008 map of parcel blocks (grain and silage maize) instead of 2015--2017 parcels, almost nothing changes (\QoneMaskPct\,\%, $p = \QoneMaskP$).
  \item \textbf{Department-wide weather.} With heat averaged over the whole department, the estimate is not significant ($p = \QoneAreaP$). \Cref{sec:maize} shows that this measure understates the 2003 heat where maize grows in departments with uplands, a pattern found after the fact; it does not explain the whole divergence, and why this measure gives a stronger association over 2010--2025 is unexplained.
  \item \textbf{Departments with uplands.} Dropping any single department leaves the estimate significant (\LooMinPct{} to \LooMaxPct\,\%), but without the six departments with uplands it falls to \NoReliefPct\,\% and is not significant ($p = \NoReliefP$). Under the rule declared before this check, the result depends on those departments, where the measure of heat also matters most; but their slope does not differ detectably from that of the others ($p = \RelInterP$), so the other departments may simply carry less information.
  \item \textbf{Another inference.} 2003 carries \LevThree\,\% of the useful variation in heat. With so much weight on one cluster, another valid inference, the cluster jackknife, gives $p = \CvThreeP$ against a Student distribution; bootstrapping the same statistic, as recommended when one cluster carries high leverage, gives $p = \CvBootP$ with a measured size of \CvBootSize\,\% (a check added later). The evidence is therefore weaker under some reasonable tests than under others.
  \item \textbf{Another crop.} Silage maize, whose irrigated and rainfed yields the statistics also report for 2010--2025, shows the same association (\FodderPct\,\% per degree-day, $p = \FodderP$, \FodderDeps{} departments), and the heat of the following summer, which cannot affect the harvest, shows none.
  \item \textbf{Heat or dry air?} When the vapour pressure deficit, the drying power of the air, enters the model with degree-days, it carries the association (\VpdJointPct\,\% per kPa, $p = \VpdJointP$) and the degree-days lose theirs. The two are too correlated to be separated, so the result is best read as concerning hot, dry summers.
  \item \textbf{Another temperature source.} With E-OBS station-based temperatures instead of SIM2, the estimate is \EobsPct\,\% ($p = \EobsP$); with satellite land-surface temperature in place of degree-days, the gap widens by \LstPct\,\% per degree Celsius ($p = \LstP$).
\end{itemize}

A pre-specified interaction with crisis-level water restrictions (2016--2025) has the opposite sign to the hypothesis: the slope of the gap on heat is not smaller where restrictions lasted longer. The other results are null, and they are reported as fully as the positive one. For soft wheat over 1959--2025, heat at the end of the cycle is associated with \QtwoPct\,\% lower yields per degree-day, but this cannot be told apart from zero ($p = \QtwoP$), and I find no sign that wheat has become more or less sensitive to heat since 1959. Out-of-sample forecasts of wheat yields from weather and satellite NDVI, built as pre-specified, are not significantly better than the simplest benchmark, the historical mean of each department.

\subsection*{What this monograph does not claim}

The maize result is an association, not the causal effect of irrigation. Irrigated and rainfed land differ in ways that may themselves interact with heat, and in dry years part of the rainfed crop sown for grain may be harvested as silage (\cref{sec:discussion}; abandoned fields stay in the statistics with a zero yield). The outcome itself is not fully measured: in the older release used for 2000--2009, the irrigated-to-rainfed ratio repeats exactly from one year to the next far more often than chance (\SaaRepOld{} of the \SaaRepTot{} repetitions of the main sample are in that release), which suggests that part of the split was carried over rather than observed (\cref{subsec:discussion-wrong}). The estimates also concern year-to-year weather shocks. They do not translate directly into the effect of a warmer climate, to which farmers would adapt.

\subsection*{Outline}

\Cref{sec:literature} reviews related work. \Cref{sec:data,sec:exposure,sec:strategy} describe the data, the measurement of heat and the empirical strategy. \Cref{sec:inference} is devoted to inference. \Cref{sec:maize,sec:wheat,sec:forecasting} report the results for maize, wheat and forecasting. \Cref{sec:discussion} discusses limitations and \cref{sec:conclusion} concludes. The appendices cover reproduction, variables, deviations from the plans, full robustness tables, the bootstrap module, the validation of the degree-day computation and further checks on maize.
